\section{Experimental design}
\label{sec:design}

\subsection{Research questions and experiment map}
The experiments address eight research questions (Table~\ref{tab:rq}). E1--E6 and E9--E11 analyse the components offline and in perception; E7 and E8 evaluate closed-loop harvesting under the pre-registered protocol of Section~\ref{sec:protocol}.

\begin{table}[t]
\centering
\caption{Research questions and experiment map. C1--C5: contributions in Section~\ref{sec:intro}.}
\label{tab:rq}
\footnotesize
\setlength{\tabcolsep}{3pt}
\begin{adjustbox}{max width=\textwidth}\begin{tabular}{>{\raggedright\arraybackslash}p{1.45cm}>{\raggedright\arraybackslash}p{4.5cm}>{\raggedright\arraybackslash}p{3.2cm}>{\raggedright\arraybackslash}p{3.4cm}>{\raggedright\arraybackslash}p{1.0cm}}
\toprule
Exp. (RQ) & Question & Data & Main metrics & Contr.\\
\midrule
E1 (RQ1) & Do the tasks require memory; how strong is the occlusion? & 640 demonstrations & Steps, picks, occlusion & C1\\
E2 (RQ2) & Does the tool-to-fruit distance locate picks; how sensitive is it to $r$, tol? & Train 200 (selection), val 199, DART-H 200 & Event P/R/F1, count exact, keyframe timing & C2\\
E3 (RQ3) & How robust are events and memory to noise, dropout and bias? & Val 199, 5 seeds per level & Event F1, count exact, keyframe recall & C2, 3\\
E4 (RQ4) & What do uniform memory and \pekm\ retain; does the coverage model hold? & Val 199--200, 5318 decisions & Event coverage, count recoverability, recency, gap & C2\\
E5 (RQ5) & Can SAM segment fruits and ripeness under occlusion and lighting shifts? & 8644 rendered images with GT masks & Instance F1, IoU, recall by visibility, latency & C3\\
E6 (RQ6) & Does RGB-D $\rightarrow$ \fdm\ $\rightarrow$ events reproduce the GT memory? & Val 199 RGB-D + proprioception & Anchor/distance error, event F1, memory agreement & C3\\
E7 (RQ7) & Does event memory improve closed-loop harvesting? & Val 50--99 $\times$ 3 seeds & Success, completion, wrong/over-pick, drop & C2, 5\\
E8 (RQ7) & Is the gain due to event information? & As E7 & Paired differences to EV & C2, 3, 5\\
E9 (RQ8) & Where do policies fail in closed loop; does DART-H help? & 10 episodes per set; E7 & TCP error by horizon; C0 vs K0 & C4\\
E10 (--) & Generalization and supporting evidence & Unseen layouts; memory benchmark & Success by quota; ablations & C2\\
E11 (RQ8) & What does the approach cost at run time? & CPU / L20 GPU & Selection and SAM latency & C2, 3\\
\bottomrule
\end{tabular}\end{adjustbox}
\end{table}

\subsection{Data and splits}
\label{sec:splits}
Training, validation and test seeds are drawn from disjoint seed ranges (base seeds 100000, 200000 and 300000), so every evaluated layout is unseen during training. All parameters of the method (Table~\ref{tab:provenance}) were fixed on training data or by physical reasoning, never on validation data. Validation episodes 0--49 were used in earlier development evaluations; the formal closed-loop evaluation (E7, E8) therefore uses validation episodes 50--99, which had not been used for any evaluation or tuning. A separate development set (episodes 200--219 of the training seed range) serves only as a go/no-go gate before the formal runs. The perception set for E5 replays validation episodes 0--49 of both tasks and renders front and wrist RGB-D images with ground-truth instance segmentation every 10th frame (4322 time steps, 8644 images); we verified that replay reproduces the recorded episodes by comparing the rendered and recorded RGB. HSV thresholds and SAM prompt settings were calibrated on 20 rendered training episodes. The offline analyses on validation demonstrations use 199 of the 200 episodes: the ground-truth file of one episode was unreadable and was regenerated afterwards from its recorded seed, with identical joint trajectories, and the coverage-by-time analysis, run after the regeneration, uses all 200.

\begin{table}[t]
\centering
\caption{Provenance of all method parameters. No parameter was tuned on validation or test data.}
\label{tab:provenance}
\footnotesize
\begin{adjustbox}{max width=\textwidth}\begin{tabular}{>{\raggedright\arraybackslash}p{4.2cm}>{\raggedright\arraybackslash}p{2.9cm}>{\raggedright\arraybackslash}p{6.0cm}}
\toprule
Parameter & Value & Source\\
\midrule
Memory budget $B$ & 16 frames & Fine-tuning configuration of the base model\\
Keyframe cap $k$ & 4, 8 & Training-set event counts (at most 3 / 5 picks per episode)\\
Event radius $r$, tolerance tol & 3 cm, 5 mm & Training-set sensitivity (E2); fruit radius 2.5--2.9 cm\\
Hysteresis exit, minimum run & $1.5r$, 3 frames & Fixed a priori\\
Coverage window & $\pm3$ frames & Fixed a priori (matching window of E2)\\
PEKM-S $\tau_d$, recent slots & 5 cm, 4 & Fixed a priori\\
DART-H $\sigma$, $\tau$, phase scales & 0.04 rad, 10, (1, 0.5, 1, 1) & Calibration on training generation: TCP displacement $\approx$2.5 cm\\
HSV thresholds & pixel-level calibration & Pixel F1 on 20 rendered training episodes\\
SAM prompt (box pad, point) & 2 px, centroid & Instance F1 on rendered training episodes\\
Vacancy / occlusion margins & 5 cm / 2 cm, 3 frames & Physical scale (fruit diameter $\approx$5 cm)\\
Grasp cue & 2 frames, within $r$ & Gripper geometry (empty-closed vs open)\\
Training (steps, batch, LR, EMA) & 12k, 64, $2.5\times10^{-5}$, 0.999 & Base-model fine-tuning recipe, unchanged\\
\bottomrule
\end{tabular}\end{adjustbox}
\end{table}

\subsection{Compared methods}
\label{sec:baselines}
\paragraph{Memory selection (E4, E7, E8)} All memories have exactly $B$ frames:
\begin{itemize}
\item \textbf{Uniform (U, K0)}: the base model's even sampling over frames 0 to $t$, identical to \pekm\ with $k=0$;
\item \textbf{Recent}: the last $B$ frames (sliding window);
\item \textbf{Random}: frame 0, frame $t$ and $B-2$ random frames (5 seeds);
\item \textbf{\pekm\ (K$k$)}: $k \in \{2, 4, 8, 12\}$ offline, $k \in \{4, 8\}$ trained for closed loop;
\item \textbf{PEKM-S}: the soft variant of Eq.~(\ref{eq:soft}), $\alpha \in \{0.25, 0.5, 1.0\}$.
\end{itemize}

\paragraph{Counterfactual controls (E8)} A change of memory content can affect a policy simply because the frames are different, whatever information they carry. To separate the two effects, each event memory (EV) is compared at test time with three controls that keep the budget, the first frame and the recent window: \textbf{EVU}, the same sampler driven by a uniform density, i.e.\ the procedure without event information; \textbf{EVR}, uniform memory in which exactly the frames that EV dropped are replaced by the same number of random non-uniform frames; and \textbf{EVA}, the anti-event selection driven by $1-w_f$ (frames far from fruits). Only EV$-$EVU and EV$-$EVR differences that favour EV, with EVA no better than U, are interpreted as an effect of event information.

\paragraph{Event detectors (E2, E3, E6)} Plain thresholding (Section~\ref{sec:events}), a causal 5-frame median filter followed by thresholding (median5), and hysteresis with a 3-frame minimum run.

\paragraph{Fruit segmentation (E5)} \textbf{HSV}, colour thresholds and connected components, a traditional baseline; \textbf{SAM-AUTO}, automatic SAM masks filtered by colour and compactness, i.e.\ zero-shot SAM without prompts (ViT-B and ViT-H, every 5th frame); \textbf{\fdm}, colour-prompted SAM with ViT-B, ViT-L and ViT-H backbones; and \textbf{SAM-GTBOX}, SAM prompted with the ground-truth box, an upper bound on mask quality given a correct prompt.

\paragraph{Perception chain (E6)}
\begin{itemize}
\item \textbf{GT}: the simulator distance, an upper bound;
\item \textbf{\fdm}: frame-0 anchors with depth and grasp evidence;
\item \textbf{\fdm\ + noisy depth}: depth noise $\sigma(z) = 1~\mathrm{mm} + 2.5~\mathrm{mm}\cdot z^2$ ($z$ in m) and 2\% missing pixels, typical of consumer RGB-D cameras;
\item \textbf{Depth-only}: no grasp evidence;
\item \textbf{Post-gust anchors}: anchors built at frame 29, after the wind gust, i.e.\ without spatial memory acquired before occlusion;
\item \textbf{Per-frame}: no anchors; the current frame's colour detections are back-projected, and an invisible fruit is unknown.
\end{itemize}

\paragraph{Training data (E9)} C0 is trained on clean demonstrations only and compared with K0, trained on clean plus DART-H demonstrations; both use uniform memory and identical settings.

\subsection{Metrics}
\label{sec:metrics}
\paragraph{Closed-loop harvesting (E7, E8)} The primary metric is the \textbf{success rate}, with the task-specific success conditions of Section~\ref{sec:tasks}. Process metrics are the \textbf{harvest completion} (fruits correctly in the basket divided by the target number), the \textbf{wrong-pick rate} (episodes with an unripe fruit in the basket, FruitOccludedRipe), the \textbf{over-pick rate} (episodes with more than $N$ fruits in the basket, FruitCount), the \textbf{drop rate} (episodes in which a fruit is dropped outside the basket) and the steps to completion and per pick. Episodes terminated by infrastructure faults are reported separately as the \textbf{engineering-error rate} and are not counted as policy failures.

\paragraph{Events (E2, E3, E6)} A detected event $[t_s, t^{\mathrm{r}}]$ is matched one-to-one, in time order, to a ground-truth detachment (the first frame at which the fruit is no longer attached) if the detachment falls within $[t_s-3, t^{\mathrm{r}}+3]$. We report precision, recall and F1; the \textbf{count exact rate}, the fraction of episodes in which the number of detected events equals the number of picks; and the contact lead and release error, in frames relative to detachment.

\paragraph{Memory (E3, E4, E6)} Memories are evaluated at every decision point (every 16 frames). A completed pick is \emph{covered} if the memory holds a frame within $\pm 3$ frames of its contact or release keyframe. We report:
\begin{itemize}
\item \textbf{event coverage}: the fraction of completed picks that are covered;
\item \textbf{run coverage}: at least one frame inside the pick's near run;
\item \textbf{count recoverability}: the fraction of decisions at which every completed pick is covered, i.e.\ the memory suffices in principle to count the picks;
\item \textbf{recency coverage}: coverage of the last, second-last and third-last pick;
\item the near-frame fraction and the median gap between memory frames, as the cost in temporal coverage;
\item for perturbed or perceived signals, the Jaccard similarity to the GT memory and the recall of GT keyframes within $\pm3$ frames.
\end{itemize}

\paragraph{Perception (E5, E6)} Segmentation is scored by instance precision, recall and F1 at IoU $\ge 0.5$, with fruits of fewer than 20 (front) or 60 (wrist) visible pixels treated as ignore regions; by the mean IoU and ripeness accuracy of matched instances; by recall stratified by the visible fraction of the fruit ($>$75\%, 25--75\%, $<$25\%); and by latency per image on an L20 GPU. Illumination robustness is measured by re-rendering every third frame photometrically as overcast or dusk, direct sun and warm low sun, with the colour thresholds kept as calibrated. Anchors are scored by recall and precision against the true fruit positions at frame 0 (5~cm gate), median and 90th-percentile localization error, ripeness accuracy, and the fraction of leaf-hidden ripe fruits that have an anchor; distances by the error of $\hat d_t$ in the near range ($d_t \le 10$~cm).

\subsection{Closed-loop evaluation protocol}
\label{sec:protocol}
The closed-loop protocol was written and sealed (SHA-256 digest recorded) before the formal runs, and the first completed run is the result, without re-runs or selection.
\begin{itemize}
\item \textbf{Episodes and seeds.} Validation episodes 50--99 of each task are evaluated with three policy sampling seeds (7, 8, 9), giving 150 trials per task and condition. A decision is made every 16 steps, and an episode ends at success, failure or 1300 steps.
\item \textbf{Conditions.} K0, K4 and K8 (uniform and \pekm, all trained on clean plus DART-H data) and C0 (uniform, clean data only).
\item \textbf{Pairing.} For each episode and seed, all conditions run in one block with the same environment seed and initial state, in randomized order, on policy servers whose identity is logged; each server holds a fixed group of conditions so that their checkpoints fit in GPU memory.
\item \textbf{Validity gates.} A run is valid only if engineering errors are at most 2\% of trials, every block is complete, and a per-decision audit confirms that the frame indices fed to the model equal those recomputed offline by the declared selection rule from the logged signal.
\item \textbf{Development gate.} Before the formal run, K0 is evaluated on the development set; the formal evaluation proceeds only if its success rate exceeds 5\% and engineering errors are at most 2 episodes, so a broken pipeline cannot consume the validation episodes.
\end{itemize}

\subsection{Statistical analysis}
\label{sec:stats}
Differences between conditions are estimated as paired differences in success rate. Confidence intervals come from a cluster bootstrap that resamples episodes and keeps all seeds and conditions of an episode together (20{,}000 resamples, fixed random seed). The primary contrasts are K4$-$K0 and K8$-$K0 per task, and the family of primary tests is corrected by the Holm method \citep{holm1979simple}. We report 95\% confidence intervals for differences, and a difference is called \emph{equivalent} if its 90\% confidence interval lies within $\pm\delta$ with $\delta=12$ percentage points (two one-sided tests, \citealp{schuirmann1987tost}). Secondary contrasts (DART-H: K0$-$C0; counterfactual controls) are reported with unadjusted intervals and labelled as secondary. Offline experiments (E2--E6) are deterministic given the data, except for noise and random baselines, which use 5 seeds; we report means over seeds.

\subsection{Hyperparameter provenance and implementation}
\label{sec:provenance}
To prevent tuning on evaluation data, every parameter was fixed before the formal evaluation, either on training data or from physical or architectural considerations (Table~\ref{tab:provenance}). The simulation uses ManiSkill3 \citep{tao2025maniskill3} with SAPIEN rendering and mplib \citep{mplib} motion planning. The policy is implemented in JAX on the openpi code base \citep{black2024pi0}, and SAM in PyTorch with the official checkpoints. Experiments ran on a server with eight NVIDIA L20 (48~GB) GPUs; in closed-loop evaluation, policy servers run on GPUs separate from the simulators.
